\chapter{लोरां श्रेणी और अवशेष प्रमेय}
\label{ch:b3:residues}

किसी \omterm{def:b3:holomorphic:holo}{समविश्लेषिक फलन} का उस बिंदु के निकट क्या होता है जहाँ वह
परिभाषित \emph{नहीं} है? उत्तर एक \omterm{def:b3:complete:complete}{पूर्ण} त्रिविभाजन है ---
विलोपनीय बिंदु, ध्रुव, अथवा सारभूत एकलता --- जिसे किसी दो-मुखी घात
श्रेणी, यानी लोरां प्रसार, में पढ़ा जाता है। उस प्रसार का एक
गुणांक, \emph{\omterm{def:b3:residues:singularities}{अवशेष}}, एकलता के चारों ओर के प्रत्येक \omterm{def:b3:holomorphic:contour}{परिरेखा समाकल}
को नियंत्रित करता है: \omterm{def:b3:residues:singularities}{अवशेष} प्रमेय कठिन निश्चित समाकलों को परिमित
बीजगणित में बदल देती है, फलनों के शून्य गिनती है (कोणांक सिद्धांत,
रूशे), और विवृत प्रतिचित्रण प्रमेय सिद्ध करती है। पहले हम कोशी की
प्रमेय को तारकाकार प्रांतों से उठाकर उसके निश्चायक, समजातता-मुक्त
रूप तक ले जाते हैं --- डिक्सन का सुघड़ तर्क --- ताकि पूरक के चारों
ओर घूर्णन संख्या शून्य वाली सभी परिरेखाएँ उपलब्ध हो जाएँ।

\section{वैश्विक कोशी प्रमेय}

\emph{चक्रक}\index{चक्रक} $\Gamma$ संवृत पथों $\gamma_1, \dots, \gamma_m$ का कोई परिमित औपचारिक
योग है; $\Gamma$ के अनुदिश समाकल और सूचकांक संगत योग होते हैं, और $\operatorname{im}\Gamma = \bigcup\operatorname{im}\gamma_j$।

\begin{theorem}[कोशी, वैश्विक रूप]
\label{thm:b3:residues:globalcauchy}
मान लीजिए $\Omega \subseteq \C$ विवृत है, $f \in \mathcal H(\Omega)$, और $\Gamma$ $\Omega$ में ऐसा चक्रक कि
\[
\operatorname{Ind}_\Gamma(w) = 0
\qquad\text{प्रत्येक } w \notin \Omega \text{ के लिए} .
\]
तब सभी $z \in \Omega\setminus\operatorname{im}\Gamma$ के लिए
\[
\frac1{2\iu\pi}\int_\Gamma\frac{f(w)}{w - z}\,\dd w =
\operatorname{Ind}_\Gamma(z)\,f(z),
\qquad\text{तथा}\qquad
\int_\Gamma f(w)\,\dd w = 0 .
\]
\index{कोशी की प्रमेय (वैश्विक)}
\end{theorem}

\begin{proof}[उपपत्ति (डिक्सन)]
$g \colon \Omega\times\Omega \to \C$ को इस प्रकार परिभाषित कीजिए
\[
g(z, w) = \begin{cases}
\dfrac{f(w) - f(z)}{w - z} & w \neq z,\\[4pt]
f'(z) & w = z .
\end{cases}
\]
\emph{$g$ \omterm{def:b3:topology:continuity}{संतत} है}: विकर्ण के बाहर यह स्पष्ट है। किसी विकर्ण
बिंदु $(a, a)$ के निकट $a$ पर $f$ का घात श्रेणी में प्रसार कीजिए
(\cref{thm:b3:holomorphic:analytic}): $f(w) - f(z) =
\sum_{n\geq1}c_n\bigl((w-a)^n - (z-a)^n\bigr)$, और प्रत्येक पद को $w - z$ से भाग देने पर ($u^n - v^n$ का
गुणनखंडन) $z, w \in D(a, r)$ के लिए
\[
g(z, w) =
\sum_{n\geq1}c_n\sum_{j=0}^{n-1}(w-a)^{\,j}(z-a)^{\,n-1-j},
\]
मिलता है, जो विकर्ण पर भी वैध है (प्रत्येक भीतरी योग $n(z-a)^{n-1}$ बन जाता
है, जिनका योग $f'(z)$ है)। छोटे $r$ के लिए यह श्रेणी $D(a,r)^2$ पर एकसमान
रूप से अभिसरित होती है ($\abs{\text{पद}} \leq n\abs{c_n}r^{n-1}$, जो त्रिज्या के भीतर योग्य है): अतः
योग \omterm{def:b3:topology:continuity}{संतत} है।

$\Omega$ पर $h(z) = \frac1{2\iu\pi}\int_\Gamma g(z, w)\,\dd w$ रखिए: यह \omterm{def:b3:topology:continuity}{संतत} है (संहतों पर $g$ का एकसमान \omterm{def:b3:topology:continuity}{सांतत्य}),
और \omterm{def:b3:holomorphic:holo}{समविश्लेषिक} भी --- मोरेरा (\cref{thm:b3:holomorphic:weierstrassconv} की कसौटी) से: किसी त्रिभुज
$T \subseteq \Omega$ के लिए फुबिनी $\int_{\partial
T}h = \frac1{2\iu\pi}\int_\Gamma\bigl(\int_{\partial T}g(z,
w)\dd z\bigr)\dd w = 0$ देता है, जिसमें भीतरी समाकल इसलिए लुप्त
होता है कि $z \mapsto g(z, w)$ $\Omega$ पर \omterm{def:b3:holomorphic:holo}{समविश्लेषिक} है ($z = w$ पर एकलता विलोपनीय
है: $g$ वहाँ \omterm{def:b3:topology:continuity}{संतत} है और अन्यत्र \omterm{def:b3:holomorphic:holo}{समविश्लेषिक} --- \cref{thm:b3:holomorphic:formula} की उपपत्ति
का विस्तार तर्क)।

\omterm{def:b3:topology:topology}{विवृत समुच्चय} $\Omega' = \{z \notin \operatorname{im}\Gamma
: \operatorname{Ind}_\Gamma(z) = 0\}$ पर $h_1(z) =
\frac1{2\iu\pi}\int_\Gamma\frac{f(w)}{w - z}\dd w$ परिभाषित कीजिए: यह $\Omega'$ पर \omterm{def:b3:holomorphic:holo}{समविश्लेषिक}
है (मोरेरा अथवा समाकल के भीतर अवकलन)। $z \in \Omega\cap\Omega'$ के लिए:
\[
h(z) = \frac1{2\iu\pi}\int_\Gamma\frac{f(w)}{w-z}\dd w
- f(z)\operatorname{Ind}_\Gamma(z) = h_1(z) .
\]
परिकल्पना से $\Omega\cup\Omega' = \C$ ($w \notin \Omega
\Rightarrow \operatorname{Ind}_\Gamma(w) = 0$), अतः $h$ और $h_1$ जुड़कर एक समग्र फलन
$H$ बना देते हैं। चूँकि $\operatorname{im}\Gamma$ के पूरक का अपरिबद्ध घटक $\Omega'$ में है
और $\abs z \to \infty$ पर $h_1(z) \to 0$ (ML परिबंध), $H$ परिबद्ध है और $0$ की ओर जाता
है: ल्यूविल (\cref{cor:b3:holomorphic:liouville}) से $H \equiv 0$। इस प्रकार $\Omega$ पर $h \equiv 0$, जो समाकल
सूत्र है। उसे स्थिर $a \in
\Omega\setminus\operatorname{im}\Gamma$ के लिए $z = a$ पर $\tilde f(w) =
(w - a)f(w)$ पर लगाने से:
\[
\frac1{2\iu\pi}\int_\Gamma f(w)\dd w =
\frac1{2\iu\pi}\int_\Gamma \frac{\tilde f(w)}{w - a}\dd w =
\operatorname{Ind}_\Gamma(a)\,\tilde f(a) = 0 .
\]
\end{proof}

\section{लोरां श्रेणी और विविक्त एकलताएँ}

\begin{theorem}[लोरां प्रसार]\label{thm:b3:residues:laurent}
मान लीजिए $f$ वलय $A = \{r < \abs{z - a} <
R\}$ ($0 \leq r < R \leq \infty$) पर \omterm{def:b3:holomorphic:holo}{समविश्लेषिक} है। तब\index{लोरां
श्रेणी}
\[
f(z) = \sum_{n\in\Z}c_n\,(z - a)^n
\qquad A \text{ पर},
\qquad
c_n = \frac1{2\iu\pi}\int_{C_\rho}\frac{f(w)}{(w -
a)^{n+1}}\,\dd w
\]
किसी भी $r < \rho < R$ के लिए (जो $\rho$ से स्वतंत्र है), और दोनों अर्ध-श्रेणी
\omterm{def:b3:topology:compact}{संहत} उपवलयों पर प्रसामान्य रूप से अभिसरित होती हैं। यह प्रसार
अद्वितीय है।
\end{theorem}

\begin{proof}
$r < \rho_1 < \abs{z - a} < \rho_2 < R$ स्थिर कीजिए और $\Gamma =
C_{\rho_2} - C_{\rho_1}$ लीजिए (बाहरी वामावर्त, भीतरी दक्षिणावर्त):
यह $A$ में ऐसा चक्रक है जिसके लिए सभी $w \notin A$ पर $\operatorname{Ind}_\Gamma(w) = 0$ (छोटे चक्र के
भीतर के बिंदु: $1 - 1 = 0$; बड़े के बाहर: $0 - 0$)। \cref{thm:b3:residues:globalcauchy} से $\operatorname{Ind}_\Gamma(z) = 1 - 0 = 1$ देता है
\[
f(z) = \frac1{2\iu\pi}\int_{C_{\rho_2}}\frac{f(w)}{w - z}\dd
w - \frac1{2\iu\pi}\int_{C_{\rho_1}}\frac{f(w)}{w - z}\dd w .
\]
पहले नाभिक का \cref{thm:b3:holomorphic:analytic} की तरह प्रसार कीजिए ($\frac{z -
a}{w - a}$ की घातें, मापांक
$< 1$): अऋणात्मक भाग $\sum_{n\geq0}c_n(z-a)^n$। दूसरे में उलटी दिशा में प्रसार कीजिए:
$\frac{-1}{w - z} = \frac1{(z-a)(1 - \frac{w - a}{z -
a})} = \sum_{m\geq0}\frac{(w-a)^m}{(z - a)^{m+1}}$, जो $C_{\rho_1}$ पर प्रसामान्य रूप से अभिसारी है: ऋणात्मक भाग $\sum_{n\leq-1}c_n(z-a)^n$,
कहे गए गुणांकों के साथ (सूचकांक $n = -m-1$)। $\rho$ से स्वतंत्रता: $C_{\rho}$
और $C_{\rho'}$ पर गुणांक समाकल $A$ में \omterm{def:b3:holomorphic:holo}{समविश्लेषिक} $\frac{f(w)}{(w-a)^{n+1}}$ के किसी
शून्य-सूचकांक चक्रक पर $\int_\Gamma$ जितने ही भिन्न होते हैं: और वह फिर
\cref{thm:b3:residues:globalcauchy} से शून्य है। अद्वितीयता: $C_\rho$ पर $\sum c_n(z-a)^n$ का $(z - a)^{-m-1}$ के सापेक्ष
पदशः समाकलन कीजिए (प्रसामान्य अभिसरण): केवल $n = m$ बचता है।
\end{proof}

\begin{definition}\label{def:b3:residues:singularities}
यदि $f$ किसी छिद्रित चक्र $D(a, R)\setminus
\{a\}$ पर \omterm{def:b3:holomorphic:holo}{समविश्लेषिक} हो, तो लोरां से
प्रसार कीजिए ($r = 0$)। तीन परस्पर अपवर्जी स्थितियाँ:\index{विविक्त
एकलता}
\begin{itemize}
  \item $n < 0$ के लिए सभी $c_n = 0$: \emph{\omterm{thm:b3:residues:riemanncw}{विलोपनीय एकलता}} (अऋणात्मक
        श्रेणी $f$ को $a$ तक \omterm{def:b3:holomorphic:holo}{समविश्लेषिक} रूप से बढ़ा देती है);
  \item परिमित संख्या के, और कम से कम एक, $n < 0$ के लिए $c_n \neq 0$:
        क्रम $m = -\min\{n : c_n \neq 0\}$ का \emph{ध्रुव}; समतुल्य रूप से $f = g/(z-a)^m$, जहाँ
        $g$ \omterm{def:b3:holomorphic:holo}{समविश्लेषिक} और $g(a)
        \neq 0$; समतुल्य रूप से $z\to
        a$ पर
        $\abs{f(z)} \to \infty$;
  \item अनंत संख्या के ऋणात्मक $c_n \neq 0$: \emph{सारभूत एकलता}।
\end{itemize}
\emph{अवशेष}\index{अवशेष} $\operatorname{Res}(f, a) = c_{-1}$ है। $\Omega$ में से ध्रुवों के किसी
समुच्चय को हटाकर \omterm{def:b3:holomorphic:holo}{समविश्लेषिक} कोई फलन $\Omega$ पर
\emph{मेरोमॉर्फिक}\index{मेरोमॉर्फिक फलन} कहलाता है।
\end{definition}

\begin{theorem}[रीमान; काज़ोराती--वाइरश्ट्रास]
\label{thm:b3:residues:riemanncw}
मान लीजिए $f$ $D(a,R)\setminus\{a\}$ पर \omterm{def:b3:holomorphic:holo}{समविश्लेषिक} है।
\begin{enumerate}
  \item (रीमान) यदि $a$ के निकट $f$ \emph{परिबद्ध} हो, तो
        एकलता विलोपनीय है।\index{विलोपनीय एकलता}
  \item (काज़ोराती--वाइरश्ट्रास) यदि $a$ सारभूत हो, तो प्रत्येक
        $\varepsilon$ के लिए $f\bigl(D(a,\varepsilon)\setminus\{a\}\bigr)$ $\C$ में सघन
        है।\index{काज़ोराती--वाइरश्ट्रास प्रमेय}
\end{enumerate}
\end{theorem}

\begin{proof}
(1) $n < 0$ और $\rho \to 0$ के लिए: $C_\rho$ पर ML से $\abs{c_n} \leq
M\rho^{-n-1}\cdot\rho\cdot\rho^{-\,n}\dots$: अतः $\abs{c_n} \leq \frac{1}{2\pi}\,2\pi\rho\cdot
M\rho^{-(n+1)} = M\rho^{-n} \to 0$
($-n > 0$ पर): इसलिए सभी ऋणात्मक गुणांक लुप्त हो जाते हैं।
(2) यदि कोई मान $b$ पास न आता हो: $a$ के निकट $\abs{f - b} \geq
\delta$, अतः
$g = 1/(f - b)$ $a$ के निकट \omterm{def:b3:holomorphic:holo}{समविश्लेषिक} और परिबद्ध है: विलोपनीय (1), और
$g$ मान $c$ के साथ बढ़ जाता है। यदि $c \neq 0$ हो, तो $f = b + 1/g$ $a$
के निकट परिबद्ध है: विलोपनीय --- जो अपवर्जित है। और यदि $c = 0$ हो,
तो $g$ का $a$ पर परिमित क्रम $m$ का शून्य है (\cref{thm:b3:holomorphic:identity}; $g \not\equiv 0$),
और $f = b + 1/g$ का क्रम $m$ का ध्रुव: फिर अपवर्जित।
\end{proof}

\section{अवशेष प्रमेय}

\begin{theorem}[अवशेष प्रमेय]\label{thm:b3:residues:residue}
मान लीजिए $\Omega$ विवृत है, $S \subseteq \Omega$ परिमित, $f \in
\mathcal H(\Omega\setminus S)$, और $\Gamma$ $\Omega\setminus S$ में
ऐसा चक्रक कि प्रत्येक $w \notin \Omega$ के लिए $\operatorname{Ind}_\Gamma(w) = 0$। तब\index{अवशेष प्रमेय}
\[
\frac{1}{2\iu\pi}\int_\Gamma f(z)\,\dd z
= \sum_{a\in S}\operatorname{Ind}_\Gamma(a)\,
\operatorname{Res}(f, a) .
\]
\end{theorem}

\begin{proof}
प्रत्येक $a \in S$ के लिए मान लीजिए $P_a(z) = \sum_{n\leq-1}c_n^{(a)}(z -
a)^n$ $a$ पर $f$ का मुख्य भाग
है: यह $\C\setminus\{a\}$ पर अभिसारी श्रेणी है ($1/(z-a)$ में उसकी त्रिज्या अनंत
है: लोरां पुच्छ सभी छोटे $\abs{z-a}$ के लिए अभिसरित होती है, अतः $(z-a)^{-1}$
की घात श्रेणी होने के कारण सर्वत्र), और वहाँ \omterm{def:b3:holomorphic:holo}{समविश्लेषिक}। तब $g = f - \sum_{a\in S}P_a$
की $S$ के प्रत्येक बिंदु पर विलोपनीय एकलताएँ हैं ($a$ पर उसके
लोरां प्रसार का कोई ऋणात्मक भाग नहीं: अन्य $P_{a'}$ $a$ पर
\omterm{def:b3:holomorphic:holo}{समविश्लेषिक} हैं), अतः $g$ $\Omega$ तक \omterm{def:b3:holomorphic:holo}{समविश्लेषिक} रूप से बढ़ जाता
है, और \cref{thm:b3:residues:globalcauchy} $\int_\Gamma g = 0$ दे देता है। अब प्रत्येक $P_a$ का समाकलन शेष है:
पदशः (\omterm{def:b3:topology:compact}{संहत} $\operatorname{im}\Gamma$ पर प्रसामान्य अभिसरण, जो $a$ से बचता है),
\[
\frac1{2\iu\pi}\int_\Gamma(z - a)^n\,\dd z = 0 \ (n \leq -2:
\text{आद्यंतर } \tfrac{(z-a)^{n+1}}{n+1}),
\qquad
\frac1{2\iu\pi}\int_\Gamma\frac{\dd z}{z - a} =
\operatorname{Ind}_\Gamma(a),
\]
अतः $\frac1{2\iu\pi}\int_\Gamma P_a =
c_{-1}^{(a)}\operatorname{Ind}_\Gamma(a)$। अब $a$ पर योग लीजिए।
\end{proof}

\begin{method}[अवशेष परिकलित करना]\label{met:b3:residues:compute}
सरल ध्रुव: $\operatorname{Res}(f, a) = \lim_{z\to a}(z -
a)f(z)$; और $g(a) \neq 0$, $h(a) = 0$, $h'(a)
\neq 0$ वाले $f = g/h$ के लिए $\operatorname{Res} = g(a)/h'(a)$।
क्रम $m$ का ध्रुव: $\operatorname{Res}(f, a) = \frac1{(m-1)!}\lim_{z\to
a}\bigl((z-a)^mf(z)\bigr)^{(m-1)}$। सारभूत एकलताएँ: प्रसार कीजिए और $c_{-1}$
पढ़ लीजिए (उदाहरणार्थ ज्ञात श्रेणियों से)। सदा सत्यापित कीजिए कि
परिरेखा वास्तव में किन ध्रुवों को घेरती है, और किस सूचकांक के
साथ।
\end{method}

\begin{example}[चार शास्त्रीय समाकल प्रकार]
\label{ex:b3:residues:integrals}
(क) \emph{$\R$ पर परिमेय}: $\int_\R\frac{\dd x}{1 +
x^4}$ के लिए ऊपरी अर्ध-समतल में किसी
बड़े अर्धवृत्त $S_R$ से बंद कीजिए: वहाँ समाकल्य $O(R^{-4})$ है, अतः
$\int_{S_R} \to 0$ (ML), और ध्रुवों $\eu^{\iu\pi/4}, \eu^{3\iu\pi/4}$ (सरल, किसी ध्रुव पर \omterm{def:b3:residues:singularities}{अवशेष} $\frac1{4z^3} = \frac{z}{4z^4} = -\frac z4$) के
साथ \omterm{def:b3:residues:singularities}{अवशेष} प्रमेय देती है
\[
\int_\R\frac{\dd x}{1 + x^4}
= 2\iu\pi\Bigl(-\frac{\eu^{\iu\pi/4}}4 -
\frac{\eu^{3\iu\pi/4}}4\Bigr)
= \frac{\pi}{\sqrt2} .
\]
(ख) \emph{फूरिये प्रकार}: $t \geq 0$ के लिए $\int_\R\frac{\eu^{\iu tx}}{1 + x^2}\dd x =
2\iu\pi\operatorname{Res}\bigl(\tfrac{\eu^{\iu tz}}{1+z^2},
\iu\bigr) = 2\iu\pi\frac{\eu^{-t}}{2\iu} = \pi\eu^{-t}$ --- ऊपरी अर्धवृत्त
इसलिए काम करता है कि वहाँ $\abs{\eu^{\iu tz}} =
\eu^{-t\operatorname{Im}z} \leq 1$; और वास्तविक भाग लेने पर: $\int_\R\frac{\cos(tx)}{1+x^2}\dd x = \pi\eu^{-\abs t}$,
जिससे \cref{exo:b3:lebesgue:10} का स्वीकृत सूत्र तय हो जाता है।
(ग) \emph{एक आवर्त पर त्रिकोणमितीय}: $z =
\eu^{\iu t}$, $\cos t = \frac{z + z^{-1}}2$, $\dd t =
\frac{\dd z}{\iu z}$
प्रतिस्थापित कीजिए: तब $\int_0^{2\pi}\frac{\dd t}{a + \cos t}$ ($a > 1$) इकाई वृत्त के भीतर अवशेषों की
गणना बन जाता है (\cref{exo:b3:residues:2})।
(घ) \emph{श्रेणी}: $f$ को $\pi\cot(\pi z)$ के साथ युग्मित कीजिए, जिसके
ध्रुव पूर्णांक हैं और \omterm{def:b3:residues:singularities}{अवशेष} $1$: सप्ताहांत समस्या $\sum n^{-2}$ और
$\sum n^{-4}$ का योग इसी प्रकार लेती है।
\end{example}

\begin{omfigure}
\begin{tikzpicture}[scale=1.15]
  \draw[->] (-2.5,0) -- (2.6,0) node[right, font=\small]
    {$\operatorname{Re}$};
  \draw[->] (0,-0.5) -- (0,2.5) node[above, font=\small]
    {$\operatorname{Im}$};
  \draw[omThm, very thick, ->] (-2,0) -- (0.3,0);
  \draw[omThm, very thick] (0.3,0) -- (2,0);
  \draw[omThm, very thick] (2,0) arc (0:90:2);
  \draw[omThm, very thick, ->] (2,0) arc (0:50:2);
  \draw[omThm, very thick] (0,2) arc (90:180:2);
  \node[omThm, font=\small] at (1.75,1.55) {$S_R$};
  \node[font=\small] at (2.2,-0.25) {$R$};
  \node[font=\small] at (-2.15,-0.25) {$-R$};
  \fill[omDef] (45:1.19) circle (1.6pt)
    node[above right, font=\small] {$\eu^{\iu\pi/4}$};
  \fill[omDef] (135:1.19) circle (1.6pt)
    node[above left, font=\small] {$\eu^{3\iu\pi/4}$};
  \fill[black!40] (-45:1.19) circle (1.6pt);
  \fill[black!40] (-135:1.19) circle (1.6pt);
\end{tikzpicture}

{\small $\int_\R\frac{\dd x}{1 +
x^4}$ के लिए अर्धवृत्त परिरेखा: $R \to \infty$ पर चाप का योगदान
$O(R^{-3})$ होता है, और \omterm{def:b3:residues:singularities}{अवशेष} प्रमेय दोनों घिरे हुए ध्रुव (नीले) गिन
लेती है। नीचे के दो ध्रुव (धूसर) बाहर हैं: सूचकांक $0$।}
\end{omfigure}

\section{कोणांक सिद्धांत और रूशे की प्रमेय}

\begin{theorem}[कोणांक सिद्धांत]
\label{thm:b3:residues:argument}
मान लीजिए $f$ $\Omega$ पर \omterm{def:b3:residues:singularities}{मेरोमॉर्फिक} है, जिसके शून्य $z_j$ (क्रम
$m_j$) और ध्रुव $p_k$ (क्रम $\mu_k$) हैं, और $\gamma$ $\Omega$ में उन
सबसे बचता हुआ कोई संवृत \omterm{def:b3:topology:pathconnected}{पथ} है, जिसके लिए $\Omega$ के बाहर $\operatorname{Ind}_\gamma = 0$। तब
\index{कोणांक सिद्धांत}
\[
\frac1{2\iu\pi}\int_\gamma\frac{f'(z)}{f(z)}\,\dd z
= \sum_j m_j\operatorname{Ind}_\gamma(z_j)
- \sum_k \mu_k\operatorname{Ind}_\gamma(p_k)
\]
(केवल परिमित संख्या के पद शून्येतर हैं)। किसी सरल वामावर्त
परिरेखा के लिए यह समाकल भीतर के शून्य घटा ध्रुव गिनता है, बहुलता
सहित --- और वह प्रतिबिंब \omterm{def:b3:topology:pathconnected}{पथ} $f\circ\gamma$ की $0$ के परितः घूर्णन संख्या
के बराबर है।
\end{theorem}

\begin{proof}
क्रम $m$ के किसी शून्य के निकट: $f = (z-a)^mg$, $g(a) \neq 0$, अतः $\frac{f'}f = \frac m{z - a} + \frac{g'}g$, जिसमें
दूसरा पद $a$ के निकट \omterm{def:b3:holomorphic:holo}{समविश्लेषिक} है: अर्थात् \omterm{def:b3:residues:singularities}{अवशेष} $m$ का एक
सरल ध्रुव। क्रम $\mu$ के किसी ध्रुव के निकट: $f = (z-a)^{-\mu}g$ \omterm{def:b3:residues:singularities}{अवशेष} $-\mu$ दे
देता है। अन्यत्र $\frac{f'}f$ \omterm{def:b3:holomorphic:holo}{समविश्लेषिक} है। (शून्येतर सूचकांक वाले
शून्य और ध्रुव $\gamma$ से घिरे किसी \omterm{def:b3:topology:compact}{संहत} क्षेत्र में होते हैं; और
तत्समता प्रमेय से वे वहाँ परिमित संख्या में हैं, $f \not\equiv 0$।) अब \cref{thm:b3:residues:residue}
लगाइए। अंतिम टिप्पणी: $\frac1{2\iu\pi}\int_\gamma\frac{f'}f =
\frac1{2\iu\pi}\int_{f\circ\gamma}\frac{\dd w}w =
\operatorname{Ind}_{f\circ\gamma}(0)$ ($w =
f(\gamma(t))$ प्रतिस्थापित कीजिए)।
\end{proof}

\begin{theorem}[रूशे]\label{thm:b3:residues:rouche}
मान लीजिए $f, g$ $\Omega$ पर \omterm{def:b3:holomorphic:holo}{समविश्लेषिक} हैं, और $\gamma$ ऐसा संवृत \omterm{def:b3:topology:pathconnected}{पथ}
जिसके लिए $\operatorname{Ind}_\gamma \in \{0,1\}$ हो और $\Omega$ के बाहर शून्य (अर्थात् एक सरल
परिरेखा)। यदि\index{रूशे की प्रमेय}
\[
\abs{g(z)} < \abs{f(z)} \qquad
\operatorname{im}\gamma \text{ पर},
\]
तो $f$ और $f + g$ के क्षेत्र $\{\operatorname{Ind}_\gamma =
1\}$ में (बहुलता सहित) शून्यों की
संख्या समान होती है।
\end{theorem}

\begin{proof}
$t \in \intcc01$ के लिए $f_t = f + tg$ का $\operatorname{im}\gamma$ पर कोई शून्य नहीं है ($\abs{f_t} \geq \abs f - \abs g >
0$), अतः
\[
N(t) = \frac1{2\iu\pi}\int_\gamma
\frac{f_t'(z)}{f_t(z)}\,\dd z
\]
सुपरिभाषित है; और वह घिरे हुए क्षेत्र में शून्य गिनता है (\cref{thm:b3:residues:argument};
कोई ध्रुव नहीं)। $N$ $t$ में \omterm{def:b3:topology:continuity}{संतत} है (समाकल्य संयुक्त रूप से
\omterm{def:b3:topology:continuity}{संतत} है, और हर नीचे से एकसमान परिबद्ध --- प्रभावी अभिसरण) और
पूर्णांक-मान वाला: अतः अचर। $N(0) = N(1)$।
\end{proof}

\begin{corollary}[विवृत प्रतिचित्रण प्रमेय]
\label{cor:b3:residues:openmapping}
किसी \omterm{def:b3:topology:connected}{संबद्ध} \omterm{def:b3:topology:topology}{विवृत समुच्चय} पर अचरेतर \omterm{def:b3:holomorphic:holo}{समविश्लेषिक फलन} विवृत
प्रतिचित्रण होता है। विशेष रूप से (फिर से) उच्चतम सिद्धांत लागू
होता है, और किसी \omterm{def:b3:holomorphic:holo}{समविश्लेषिक} एकैकी आच्छादक प्रतिचित्रण का प्रतिलोम
\omterm{def:b3:holomorphic:holo}{समविश्लेषिक} होता है।\index{विवृत प्रतिचित्रण प्रमेय (समविश्लेषिक)}
\end{corollary}

\begin{proof}
मान लीजिए $f(a) = b$; $f - b$ का $a$ पर किसी परिमित क्रम $m
\geq 1$ का शून्य
है (तत्समता प्रमेय: $f \not\equiv b$)। ऐसा $r$ चुनिए कि $\bar D(a, r)\setminus\{a\}$ पर $f - b$
शून्य-रहित हो (\omterm{thm:b3:holomorphic:identity}{विविक्त शून्य}) और मान लीजिए $\delta = \min_{\abs{z - a} =
r}\abs{f(z) - b} > 0$। $\abs{w - b} < \delta$ के लिए:
वृत्त पर $\abs{(b - w)} < \delta \leq \abs{f - b}$, अतः रूशे ($f - b$ बनाम अचर $b - w$) कहती है कि $f - w$
के $D(a, r)$ में ठीक $m$ शून्य हैं: अर्थात् ऐसा प्रत्येक $w$
प्राप्त होता है --- $f(D(a,r)) \supseteq D(b, \delta)$: विवृत। उच्चतम सिद्धांत: अचरेतर $f$
के लिए $\abs f$ का कोई आंतरिक अधिकतम असंभव है, क्योंकि $f(a)$ के आसपास
उसके प्रतिबिंब में बड़े मापांक के बिंदु होते हैं। प्रतिलोम: कोई
\omterm{def:b3:holomorphic:holo}{समविश्लेषिक} एकैकी आच्छादक $f$ विवृत है, अतः $f^{-1}$ \omterm{def:b3:topology:continuity}{संतत} है;
$a$ पर $f - f(a)$ का शून्य सरल है ($m \geq 2$ होने पर निकटवर्ती मानों के
$m$ पूर्वप्रतिबिंब मिलते --- और वे भिन्न होते, क्योंकि $f'$
केवल विविक्त बिंदुओं पर लुप्त होता है, अतः $a$ के निकट प्ररूपी
छोटे $w$ के लिए $f - w$ के $m$ शून्य सरल और भिन्न होते: जो
एकैकीपन का खंडन है); तब $f'(a) \neq 0$ और $f^{-1}$ का अंतर भागफल अभिसरित होता
है: $\bigl(f^{-1}\bigr)'(b) =
1/f'(a)$।
\end{proof}

\section{अभ्यास}

\begin{exercise}[$\star$]\label{exo:b3:residues:1}
$0$ पर एकलता का वर्गीकरण कीजिए और \omterm{def:b3:residues:singularities}{अवशेष} परिकलित कीजिए:
\[
\frac{\sin z}{z},\qquad
\frac{\eu^z - 1}{z^2},\qquad
\frac{1}{z(z-1)^2},\qquad
\frac{\cos z}{z^3},\qquad
\eu^{1/z},\qquad
\frac1{\sin z} .
\]
तीसरे का $z = 1$ पर और अंतिम का $z = \pi$ पर \omterm{def:b3:residues:singularities}{अवशेष} भी दीजिए।
\end{exercise}

\begin{exercise}[$\star$]\label{exo:b3:residues:2}
$a > 1$ के लिए $z = \eu^{\iu t}$ के माध्यम से परिकलित कीजिए:
\[
\int_0^{2\pi}\frac{\dd t}{a + \cos t}
= \frac{2\pi}{\sqrt{a^2 - 1}} .
\]
सीमांत व्यवहार $a \to 1^+$ और $a \to \infty$ जाँचिए।
\end{exercise}

\begin{exercise}[$\star\star$]\label{exo:b3:residues:3}
अर्धवृत्त परिरेखाओं से परिकलित कीजिए, और चाप आकलन उचित ठहराइए:
\[
\int_\R\frac{x^2}{1 + x^6}\,\dd x = \frac\pi3,
\qquad
\int_\R\frac{\dd x}{(1 + x^2)^{2}} = \frac\pi2
\quad\text{(\cref{exo:b3:fouriertransform:3} की पुनर्व्युत्पत्ति)}.
\]
\end{exercise}

\begin{exercise}[$\star\star$]\label{exo:b3:residues:4}
$t \geq 0$ और $a > 0$ के लिए सिद्ध कीजिए:
\[
\int_\R\frac{\cos(tx)}{x^2 + a^2}\,\dd x =
\frac{\pi}{a}\,\eu^{-at},
\]
और प्रतिलोमन से $x \mapsto \eu^{-a\abs
x}$ का फूरिये रूपांतर निकालिए --- \cref{exo:b3:fouriertransform:1} से
तुलना करते हुए।
\end{exercise}

\begin{exercise}[$\star\star$]\label{exo:b3:residues:5}
$f(z) = \dfrac1{(z-1)(z-2)}$ का तीनों क्षेत्रों $\abs z < 1$, $1 < \abs z < 2$, $\abs z >
2$ में \omterm{thm:b3:residues:laurent}{लोरां श्रेणी} में
प्रसार कीजिए। तीनों प्रसार भिन्न क्यों हैं? समझाइए कि दूसरे और
तीसरे प्रसार में $z^{-1}$ का गुणांक $0$ पर $f$ का \omterm{def:b3:residues:singularities}{अवशेष}
\emph{नहीं} है ($f$ की वहाँ कोई एकलता ही नहीं है), और $f$ के
वास्तविक \omterm{def:b3:residues:singularities}{अवशेष} $1$ पर तथा $2$ पर परिकलित कीजिए।
\end{exercise}

\begin{exercise}[$\star\star$]\label{exo:b3:residues:6}
(क) दर्शाइए कि $\eu^{1/z}$ की $0$ पर सारभूत एकलता है और
काज़ोराती--वाइरश्ट्रास हाथ से सत्यापित कीजिए: किसी भी $w \neq 0$ के लिए
$\eu^{1/z} =
w$ स्पष्ट रूप से हल कीजिए और $0$ के मनमाने निकट हल प्रस्तुत
कीजिए।
(ख) दर्शाइए कि $\abs{\eu^{1/z}}$ $0$ के प्रत्येक छिद्रित \omterm{def:b3:topology:topology}{प्रतिवेश} पर
अपरिबद्ध है, फिर भी $\eu^{1/z}$ का कोई ध्रुव नहीं है: कौन-सी सीमा विफल
होती है?
\end{exercise}

\begin{exercise}[$\star\star$]\label{exo:b3:residues:7}
रूशे से गिनिए:
(क) $\abs z < 1$ में $z^7 - 4z^3 + z - 1$ के शून्य;
(ख) $\abs z < 1$ में और $1 <
\abs z < 2$ में $z^4 + 5z + 1$ के शून्य;
(ग) दालांबेर--गाउस पुनः सिद्ध कीजिए: घात $n$ के किसी मोनिक
बहुपद के किसी बड़े चक्र में $n$ शून्य होते हैं \emph{($z^n$ से
तुलना कीजिए)}।
\end{exercise}

\begin{exercise}[$\star\star\star$]\label{exo:b3:residues:8}
(हुर्विट्स) मान लीजिए संहतों पर एकसमान रूप से $f_n \to f$, जहाँ $f_n \in
\mathcal H(\Omega)$,
$\Omega$ \omterm{def:b3:topology:connected}{संबद्ध}, $f \not\equiv 0$।
(क) दर्शाइए कि यदि सभी $f_n$ शून्य-रहित हों, तो $f$ भी।
\emph{(यदि $f(a) = 0$: $a$ के चारों ओर किसी छोटे वृत्त पर कोणांक
सिद्धांत, और $\int f_n'/f_n$ में सीमा लेने के लिए \cref{thm:b3:holomorphic:weierstrassconv}।)}
(ख) दर्शाइए कि यदि सभी $f_n$ एकैकी हों, तो $f$ एकैकी है अथवा
अचर। \emph{($\Omega\setminus\{w\}$ पर $z \mapsto f_n(z) - f_n(w)$ पर (क) लगाइए।)}
\end{exercise}

\begin{exercise}[$\star\star\star$]\label{exo:b3:residues:9}
$n \geq 2$ के लिए त्रिज्यखंड $\{0 \leq \arg z \leq \frac{2\pi}n,\ \abs z
\leq R\}$ की परिसीमा पर $\frac1{1 + z^n}$ का समाकलन कीजिए
और निकालिए
\[
\int_0^{+\infty}\frac{\dd x}{1 + x^n} =
\frac{\pi}{n\,\sin(\pi/n)} .
\]
$n = 2$ को $\arctan$ के सापेक्ष जाँचिए, और सीमा $n \to
\infty$ भी।
\end{exercise}

\begin{exercise}[$\star\star$]\label{exo:b3:residues:10}
मान लीजिए $f$ $\deg(\text{हर})
\geq \deg(\text{अंश}) + 2$ वाला परिमेय फलन है। दर्शाइए कि $f$ के
\emph{समस्त} अवशेषों का योग शून्य है \emph{(बड़े से बड़े वृत्तों
पर समाकलन कीजिए)}। इसका उपयोग करके $\frac1{z(z-1)(z-2)}$ का आंशिक भिन्न विघटन बिना
किसी रैखिक बीजगणित के पुनः परिकलित कीजिए।
\end{exercise}

\begin{exercise}[$\star\star\star$]\label{exo:b3:residues:11}
(कुंजीछिद्र: ऑयलर का परावर्तन समाकल) $0 < a < 1$ के लिए
\[
I(a) = \int_0^{\infty}\frac{x^{a-1}}{1 + x}\,\dd x =
\frac{\pi}{\sin(\pi a)}
\]
परिकलित कीजिए, और वह भी $f(z) = \frac{z^{a-1}}{1+z} =
\frac{\eu^{(a-1)\log z}}{1 + z}$ (लघुगणक $\R_+$ के अनुदिश कटा, $\arg z \in \intoo0{2\pi}$)
का कुंजीछिद्र परिरेखा पर समाकलन करके: कट के ऊपर से $\varepsilon$ से $R$
तक बाहर, $C_R$ के चारों ओर, कट के नीचे से वापस, और $C_\varepsilon$ के चारों
ओर। उचित ठहराइए: दोनों सीधे खंड गुणक $\eu^{2\iu\pi(a-1)}$ जितने भिन्न होते हैं,
वृत्तों के योगदान लुप्त हो जाते हैं ($R^{a-1}\cdot R \to 0$ और $\varepsilon^{a-1}\cdot\varepsilon \to 0$), और एकमात्र
ध्रुव $z = -1$ का \omterm{def:b3:residues:singularities}{अवशेष} $\eu^{\iu\pi(a - 1)}$ है। $\Gamma(a)\Gamma(1 - a) = \frac\pi{\sin\pi a}$ भी निकालिए \emph{(\cref{pb:b3:lebesgue:1} से
$\Gamma(a)\Gamma(1-a) = B(a, 1-a)$ लिखिए और $t =
\frac{x}{1+x}$ प्रतिस्थापित कीजिए)}।
\end{exercise}

\begin{exercise}[$\star\star$]\label{exo:b3:residues:12}
(कोणांक सिद्धांत से शून्य गिनना, संख्यात्मक रूप से) मान लीजिए
$P(z) = z^4 + 8z + 1$।
(क) इकाई चक्र में कितने शून्य हैं? \emph{($8z + 1$ के सापेक्ष रूशे।)}
(ख) वलय $1 < \abs z < 3$ में कितने? \emph{($\abs z = 3$ पर $z^4$ के सापेक्ष रूशे।)}
तीक्ष्ण कीजिए: दर्शाइए कि प्रत्येक शून्य का मापांक $< 2.1$ है।
(ग) दक्षिण अर्ध-समतल में कितने? \emph{(पहले $\abs z = 2$ पर गिनिए; फिर
काल्पनिक अक्ष के प्रतिबिंब का अनुसरण कीजिए: $P(\iu t) = t^4 + 1 + 8\iu t$ का वास्तविक भाग
सर्वत्र धनात्मक है, अतः अक्ष पर कोई शून्य नहीं, और उसके अनुदिश
कोणांक का विचरण परिकलनीय है --- किसी बड़े अर्ध-चक्र से निष्कर्ष
निकालिए।)}
\end{exercise}

\section{समस्या: कोटिस्पर्शज्या से
\texorpdfstring{$\zeta(2k)$}{ZETA(2k)}}

\begin{problem}[{सप्ताहांत समस्या --- अवशेषों से $\sum n^{-2k}$ का योग}]\label{pb:b3:residues:1}
\omterm{def:b3:residues:singularities}{अवशेष} प्रमेय श्रेणियों का योग ले लेती है: किसी परिमेय फलन को
$\pi\cot(\pi z)$ के साथ युग्मित करने पर, जिसके ध्रुव पूर्णांकों पर बैठते हैं,
$\sum_{n}f(n)$ अवशेषों की गणना बन जाता है। हम यह विधि सिद्ध करते हैं और
$\zeta(2) = \frac{\pi^2}{6}$ तथा $\zeta(4) =
\frac{\pi^4}{90}$ परिकलित करते हैं --- वही मान जो दूसरे वर्ष में
फूरिये श्रेणी से और \cref{ch:b3:spectral} में संकारक अनुरेखों से मिले थे, अब
परिरेखा समाकलन से।

\medskip
\noindent\textbf{भाग I --- कोटिस्पर्शज्या नाभिक।}
\begin{enumerate}
  \item दर्शाइए कि $\pi\cot(\pi z)$ $\C$ पर \omterm{def:b3:residues:singularities}{मेरोमॉर्फिक} है, जिसके सरल ध्रुव
        ठीक $z = n \in \Z$ पर हैं और प्रत्येक का \omterm{def:b3:residues:singularities}{अवशेष} $1$ है
        \emph{($\lim_{z\to n}(z-n)\pi\cot\pi z$ परिकलित कीजिए)}।
  \item $0$ पर लोरां प्रसार का आरंभ परिकलित कीजिए:
        \[
        \pi\cot(\pi z) = \frac1z - \frac{\pi^2}{3}\,z -
        \frac{\pi^4}{45}\,z^3 + O(z^5) ,
        \]
        और वह भी $\cos$ की घात श्रेणी को $\sin$ की घात श्रेणी से
        भाग देकर (भाग को उचित ठहराइए: $\frac{\sin\pi z}{\pi z}$ $0$ के निकट
        \omterm{def:b3:holomorphic:holo}{समविश्लेषिक} और शून्येतर है, अतः उसका व्युत्क्रम
        \omterm{def:b3:holomorphic:holo}{समविश्लेषिक} है; कोटि $3$ तक गुणांक पहचानिए)।
  \item मान लीजिए $C_N$ शीर्षों $(\pm1\pm\iu)(N + \frac12)$ वाले वर्ग की परिसीमा है।
        दर्शाइए कि प्रत्येक $N \geq
        1$ के लिए $C_N$ पर $\abs{\cot(\pi z)} \leq 2$।
        \emph{(ऊर्ध्वाधर भुजाओं पर $\cot(\pi(\pm(N +
        \frac12) + \iu y)) = \mp\tan(\iu\pi y)$, जिसका मापांक $\abs{\tanh(\pi y)} \leq 1$ है;
        और क्षैतिज भुजाओं $\abs y = N + \frac12$ पर $\abs{\cot(\pi(x\pm\iu y))} \leq \coth(\pi y) \leq
        \coth(\pi/2) < 1.1$ को परिबद्ध कीजिए।)}
\end{enumerate}

\medskip
\noindent\textbf{भाग II --- योग प्रमेय।}
\begin{enumerate}[resume]
  \item मान लीजिए $f$ परिमेय है, पूर्णांकों पर \omterm{def:b3:holomorphic:holo}{समविश्लेषिक}, और
        $\deg(\text{हर}) \geq \deg(\text{अंश}) + 2$। $C_N$ पर \omterm{def:b3:residues:singularities}{अवशेष} प्रमेय तथा प्रश्न 3 के परिबंध का
        उपयोग करके सिद्ध कीजिए:
        \[
        \lim_{N\to\infty}\ \sum_{n = -N}^{N} f(n)
        = -\sum_{p\ \text{जो } f \text{ का ध्रुव है}}
        \operatorname{Res}\bigl(\pi\cot(\pi z)f(z),\,p\bigr).
        \]
  \item तर्क को घात शर्त की आवश्यकता कहाँ पड़ती है? उदाहरण से
        दर्शाइए ($f(z) = 1/(z + \frac12)$ लीजिए) कि धीमे क्षय के लिए सममित सीमा फिर
        भी विद्यमान हो सकती है जबकि दो-मुखी श्रेणी अपसरित होती है
        --- और तब यह सूत्र \emph{मुख्य मान} परिकलित करता है।
\end{enumerate}

\medskip
\noindent\textbf{भाग III --- मान।}
\begin{enumerate}[resume]
  \item इस विधि को $f(z) = 1/z^2$ पर लगाइए: यहाँ $f$ का ध्रुव किसी
        पूर्णांक \emph{पर} है, अतः तर्क सीधे चलाइए --- $C_N$ पर
        $g(z) = \frac{\pi\cot(\pi
        z)}{z^2}$ का समाकलन कीजिए, दर्शाइए कि समाकल $\to 0$, और प्रश्न 2
        से $\operatorname{Res}(g, 0)$ परिकलित कीजिए। निष्कर्ष निकालिए:
        \[
        2\sum_{n\geq1}\frac1{n^2} = \frac{\pi^2}{3},
        \qquad \zeta(2) = \frac{\pi^2}6 .
        \]
  \item $g(z) = \frac{\pi\cot(\pi z)}{z^4}$ के साथ वही कीजिए: $\operatorname{Res}(g, 0)$ परिकलित कीजिए और $\zeta(4) = \frac{\pi^4}{90}$
        निकालिए।
  \item व्यापक प्रतिरूप समझाइए: प्रत्येक $k \geq 1$ के लिए $\zeta(2k)$
        $0$ पर $\pi\cot(\pi
        z)$ के लोरां प्रसार में $z^{2k-1}$ के गुणांक का
        $-\frac12$ गुना है --- अर्थात् $\pi^{2k}$ का कोई परिमेय गुणज। प्रश्न
        2 के भाग को एक चरण आगे बढ़ाकर $\zeta(6)$ परिकलित कीजिए। यह विधि
        $\zeta(3)$ के बारे में क्या कहती है --- और कुछ क्यों नहीं
        कहती?
\end{enumerate}

\medskip
\noindent\textbf{भाग IV --- कोटिस्पर्शज्या का आंशिक भिन्न
प्रसार।}
\begin{enumerate}[resume]
  \item $w \in \C\setminus\Z$ स्थिर कीजिए और भाग II की विधि $f(z) = \dfrac{1}{(z - w)(z + w)}$ पर लगाइए ---
        ध्यान दीजिए कि अब $\pi\cot(\pi z)f(z)$ के $\pm w$ पर अतिरिक्त सरल ध्रुव हैं,
        जिनके \omterm{def:b3:residues:singularities}{अवशेष} भी गणना में जुड़ने चाहिए। आंशिक भिन्न प्रसार
        \[
        \pi\cot(\pi w) = \frac1w +
        \sum_{n\geq1}\frac{2w}{w^2 - n^2},
        \]
        निकालिए, जिसमें श्रेणी $\C\setminus\Z$ के \omterm{def:b3:topology:compact}{संहत} उपसमुच्चयों पर
        प्रसामान्य रूप से अभिसरित होती है।
  \item इस प्रसार से, प्रत्येक पद को $w$ की घातों में फैलाकर
        (अदला-बदली उचित ठहराइए), प्रश्न 2 जैसे \emph{वही} लोरां
        गुणांक पुनः प्राप्त कीजिए --- वृत्त पूरा हो जाता है:
        ऑयलर का सूत्र $\sum\frac1{n^2} = \frac{\pi^2}6$ कोटिस्पर्शज्या के दोनों चेहरों में
        $w$ का गुणांक है। फूरिये-श्रेणी उपपत्ति (दूसरा वर्ष) और
        अनुरेख उपपत्ति (\cref{pb:b3:spectral:1}) से तुलना कीजिए: तीन सिद्धांत, एक
        संख्या।
\end{enumerate}

\medskip
\noindent\textbf{भाग V --- ज्या के लिए ऑयलर का गुणनफल।} प्रश्न 9
का प्रसार किसी अनंत गुणनफल का लघुगणकीय अवकलज है; अब हम ऑयलर के
1734 के गुणनखंडन को ईमानदारी से सिद्ध करते हैं।
\begin{enumerate}[resume]
  \item $N \geq 1$ के लिए $P_N(z) =
        z\prod_{n=1}^{N}\bigl(1 - \frac{z^2}{n^2}\bigr)$ रखिए। दर्शाइए कि $P_N$ प्रत्येक चक्र
        $\bar D(0, R)$ पर एकसमान रूप से किसी समग्र फलन $P$ पर अभिसरित
        होता है, जिसके शून्य ठीक पूर्णांक हैं और सभी सरल।
        \emph{($n
        \geq 2R$ के लिए गुणनखंड को \cref{exo:b3:holomorphic:3} के मुख्य लघुगणक के
        साथ $\exp\log(1 -
        z^2/n^2)$ के रूप में लिखिए, $\abs u \leq \frac12$ के लिए श्रेणी के
        माध्यम से $\abs{\log(1+u)}
        \leq 2\abs u$ परिबद्ध कीजिए, और लघुगणकों के प्रसामान्य
        अभिसारी योग का चरघातांक लीजिए; शेष परिमित गुणनखंड एक
        बहुपद हैं। \cref{thm:b3:holomorphic:weierstrassconv} से निष्कर्ष निकालिए।)}
  \item दर्शाइए कि $\C\setminus\Z$ पर
        \[
        \frac{P'(z)}{P(z)} = \frac1z +
        \sum_{n\geq1}\frac{2z}{z^2 - n^2} = \pi\cot(\pi z)
        \]
        \emph{(परिमित गुणनफलों का अवकलन कीजिए, \cref{thm:b3:holomorphic:weierstrassconv} और $\Z$ के
        बाहर $P$ की शून्य-रहितता का उपयोग करते हुए सीमा लीजिए,
        और प्रश्न 9 उद्धृत कीजिए)}।
  \item दर्शाइए कि $Q = \sin(\pi z)/P(z)$ $Q' = 0$ वाले किसी शून्य-रहित समग्र फलन तक
        बढ़ जाता है, और \emph{ऑयलर का गुणनफल} निष्कर्ष के रूप में
        निकालिए:
        \[
        \sin(\pi z) = \pi z\prod_{n\geq1}
        \Bigl(1 - \frac{z^2}{n^2}\Bigr)
        \qquad (z \in \C) .
        \]
  \item (वालिस, 1655) $z = \frac12$ पर मान निकालिए:
        \[
        \frac\pi2 = \prod_{n\geq1}\frac{4n^2}{4n^2 - 1}
        = \lim_{N\to\infty}
        \frac{2\cdot2\cdot4\cdot4\cdots(2N)(2N)}
        {1\cdot3\cdot3\cdot5\cdots(2N-1)(2N+1)} .
        \]
  \item $\abs z < 1$ के लिए गुणनफल के लघुगणक का द्विक श्रेणी में प्रसार
        कीजिए (पुनर्विन्यास उचित ठहराइए) और $\sin(\pi z) = \pi z - \frac{\pi^3}6z^3 + \cdots$ में $z^3$ के
        गुणांक का मिलान करके $\zeta(2) =
        \frac{\pi^2}6$ पुनः प्राप्त कीजिए --- यही
        ऑयलर की संख्या का गुणनफल वाला चेहरा है।
\end{enumerate}

\medskip
\noindent\textbf{भाग VI --- सहोदर नाभिक।} कोटिस्पर्शज्या के
सहोदर हैं; प्रत्येक अपने ही श्रेणी-कुल का मोल तय करता है।
\begin{enumerate}[resume]
  \item प्रश्न 9 के प्रसार का पदशः अवकलन कीजिए (\cref{thm:b3:holomorphic:weierstrassconv} से उचित)
        और $\C\setminus\Z$ के संहतों पर प्रसामान्य रूप से
        \[
        \frac{\pi^2}{\sin^2(\pi z)} =
        \sum_{n\in\Z}\frac1{(z - n)^2} .
        \]
        प्राप्त कीजिए।
  \item $z = \frac12$ पर मान निकालिए: $\sum_{m\geq0}\frac1{(2m+1)^2} = \frac{\pi^2}8$; और पूर्णांकों को सम-विषमता
        से बाँटकर $\zeta(2)$ एक बार फिर पुनः प्राप्त कीजिए।
  \item द्विगुणन सर्वसमिका $\tan\theta =
        \cot\theta - 2\cot(2\theta)$ सत्यापित कीजिए और $\frac12 + \Z$ से बचने
        वाले संहतों पर प्रसामान्य रूप से
        \[
        \pi\tan(\pi z) =
        \sum_{m\geq0}\frac{8z}{(2m+1)^2 - 4z^2} ,
        \]
        निकालिए।
  \item $0$ के परितः प्रसार कीजिए ($\abs z < \frac12$; फिर फुबिनी): $\lambda(s) =
        \sum_{m\geq0}(2m+1)^{-s}$
        के साथ
        \[
        \pi\tan(\pi z) =
        \sum_{k\geq0}8\cdot4^k\,\lambda(2k+2)\,z^{2k+1} ;
        \]
        $\tan u = u + \frac{u^3}3 + O(u^5)$ से तुलना करके $\lambda(2) = \frac{\pi^2}8$ पुनः प्राप्त कीजिए और $\lambda(4) = \frac{\pi^4}{96}$ भी;
        फिर $\lambda(4) = (1 -
        2^{-4})\,\zeta(4)$ के माध्यम से $\zeta(4) = \frac{\pi^4}{90}$ की जाँच कीजिए।
  \item $\frac1{\sin\theta} = \cot\frac\theta2 -
        \cot\theta$ सत्यापित कीजिए और निकालिए
        \[
        \frac{\pi}{\sin(\pi z)} = \frac1z +
        \sum_{n\geq1}(-1)^n\,\frac{2z}{z^2 - n^2} .
        \]
        चिह्नों को पूर्णांकों पर $\pi/\sin(\pi z)$ के अवशेषों के सापेक्ष
        जाँचिए।
  \item $z$ का गुणांक पढ़ लीजिए: $\eta(2) =
        \sum_{n\geq1}\frac{(-1)^{n-1}}{n^2} =
        \frac{\pi^2}{12}$, और संगति $\eta(2) = (1 - 2^{1-2})\,\zeta(2)$ की पुष्टि
        कीजिए।
  \item (समापन) प्रश्न 20 के प्रसार का $z =
        \frac12$ पर मान निकालिए और
        \emph{लाइब्नित्स का सूत्र}
        \[
        \frac\pi4 = 1 - \frac13 + \frac15 - \frac17 +
        \cdots
        \]
        निकालिए। एक छोटे अनुच्छेद से समाप्त कीजिए: प्रत्येक
        अंकगणित के लिए एक नाभिक --- कौन-सा नाभिक किस श्रेणी-कुल
        का मोल तय करता है, और वे सब $\zeta(3)$ के प्रति संरचनात्मक रूप
        से अंधे क्यों हैं।

\end{enumerate}
\medskip
\noindent\textbf{भाग VII --- पूरी मूल्य सूची: बर्नूली
संख्याएँ।}
\begin{enumerate}[resume]
  \item $\pi
        z\cot(\pi z)$ के आंशिक भिन्न प्रसार को बर्नूली संख्याओं के जनक
        फलन ($\frac{w}{\eu^w - 1} =
        \sum_n\frac{B_n}{n!}w^n$, \cref{pb:b3:holomorphic:1}, भाग VI) के साथ मिलाइए: \[
        \pi z\cot(\pi z) = \iu\pi z +
        \frac{2\iu\pi z}{\eu^{2\iu\pi z} - 1}
        \] से (पहले
        यह सर्वसमिका सिद्ध कीजिए) संवृत रूप
        \[
        \zeta(2k) = (-1)^{k+1}\,
        \frac{(2\pi)^{2k}\,B_{2k}}{2\,(2k)!}
        \qquad (k \geq 1).
        \]
        निकालिए।
  \item इस सूत्र को $B_2 = \frac16$, $B_4 =
        -\frac1{30}$, $B_6 = \frac1{42}$ के सापेक्ष सत्यापित कीजिए:
        $\zeta(2) =
        \frac{\pi^2}6$, $\zeta(4) = \frac{\pi^4}{90}$ पुनः प्राप्त कीजिए, और $\zeta(6) = \frac{\pi^6}{945}$ परिकलित कीजिए।
  \item (ऑयलर का पुनरावर्तन) $\bigl(z\cot z\bigr)' = \cot z - z(1 + \cot^2z)$ के दोनों पक्षों का प्रसार
        कीजिए --- अथवा कोटिस्पर्शज्या श्रेणी का सीधे वर्ग कीजिए ---
        और सिद्ध कीजिए
        \[
        \Bigl(k + \frac12\Bigr)\zeta(2k) =
        \sum_{j=1}^{k-1}\zeta(2j)\,\zeta(2k - 2j)
        \qquad (k \geq 2),
        \]
        और जाँचिए कि यह $\zeta(2)$ से $\zeta(4)$ और $\zeta(2), \zeta(4)$ से $\zeta(6)$ परिकलित
        कर देता है --- अर्थात् एकमात्र बीज $\frac{\pi^2}6$ से समस्त सम ज़ीटा
        मान, और कोई नया समाकलन नहीं।
\end{enumerate}
\end{problem}
